\section{Audit findings and proposed manuscript corrections}
\label{audit:sec:corrections}

The research does not require withdrawing the source's global zero
bound, its eventual uniform saturation theorem, its $n=1$ monotonicity
theorem, or its all-conductor symbol kernel theorem. Those results
survive the present analysis. The following corrections and status
updates are narrower and should be integrated with their supporting
proofs.

\subsection{A logical wording correction in the rank conjecture}

In \path{chapters/05-signed-kernels.tex}, Conjecture
\texttt{signed:conj:rank} states two absolute ranks and then describes
their difference as an equivalent statement. The difference follows
from the two absolute formulas, but it does not recover them without
another premise. Replace ``Equivalently'' by ``Consequently'' and
``The equivalence'' by ``This implication.'' This is a logical wording
correction, independent of whether the conjecture is true. Theorem
\ref{fg:thm:rank} now proves both absolute ranks, so after its proof
has been integrated the conjecture can be promoted to a theorem.

The attached patch makes only those wording changes and the duplicated
phrase correction below. It deliberately leaves the conjecture
environment intact: installing a new theorem label without its proof
and dependencies would be a poor intermediate repository state.

\subsection{Boundary regularization must remain explicit}

The removed coordinate $X_1^{0,1}$ and the retained difference row are
essential parts of the matrix definition. A computation that discards
both boundary rows, or retains them separately as homogeneous rows,
solves a different problem. The canonical defect $cY^n$ in Lemma
\ref{fg:lem:boundary} identifies exactly what the valid regularization
preserves. The matrix generator includes column order, row labels, and
the weight-five matrix itself so future implementations can compare
the actual presentation, not merely a rank value.

Similarly, the new universal splitting concerns only the specified
single-product coefficient system. Adding distribution or higher-depth
relations can change the residual kernel. The theorem should not be
paraphrased as a numerical dimension conjecture for all cyclotomic
multiple polylogarithms.

\subsection{The angular expansion is no longer an open construction}

The paragraph following the source's angular numerical table describes
a systematic expansion beyond four terms as a separate research
problem. Theorem~\ref{angular:thm:allorders} resolves that problem at
every fixed exponential accuracy, and extends the uniformity from
$b\geq1$ to $b>0$. The correct replacement should also mention
Theorem~\ref{angular:thm:divergence}: the infinite grouped series has
unbounded terms at fixed $a,b$. Ordinary convergence is therefore
ruled out, not merely unproved.

An analytic implicit-function argument on the entire complex Fourier
series would be invalid without additional work: its frequencies grow
exponentially off the real axis. The finite-mode argument in this
article is part of the proof, not an implementation detail that may
be dropped from an abbreviated presentation.

\subsection{Higher-index monotonicity needs counterexamples, not extrapolation}

The source proves a decreasing branch at $n=1$. It does not thereby
prove decreasing branches at larger indices. Theorem
\ref{lerchboundary:thm:counterexample} gives a strictly positive initial
velocity for the smallest $n=4,k=3$ branch. Theorem
\ref{lerchboundary:thm:n2global} gives an entire $n=2,k=1$ branch that
is nonmonotone on $[0,1]$. These are useful additions to the discussion
of higher-index motion; they do not contradict the correctly restricted
$n=1$ theorem.

The regularity statement must distinguish the two endpoints. At a
simple elementary zero, analytic dependence on $\rho$ extends through
$0$. At a simple Stieltjes endpoint zero, Theorem
\ref{lerchboundary:thm:abel-root} gives exactly $C^{k-1}$ regularity,
with a divergent $k$th derivative. An unqualified claim of analyticity
on a neighborhood of the closed parameter interval would be false.

In the proof of the source result \texttt{half:zero:first}, replace
``The endpoint signs in the endpoint signs and the Laplace representation
above'' by ``The endpoint signs and the Laplace representation above.''
This is an editorial duplication only.

\subsection{Formal Herglotz obstructions and analytic evaluations}

The general companion classification must retain the condition modulo
$2e$. For example, the pair $e=8,o=3$ passes the corresponding tests
modulo $o$ and $e$, but fails $o^2\equiv1\pmod{2e}$. A criterion that
simply reuses the two moduli from the source's $F$-classification would
give an incorrect answer for $J(3/8)$.

The term ``no reduction'' in this setting means no reduction of the
specified invariant pre-Bloch element. The boundary argument is not
a transcendence theorem for actual values of $J$. Conversely, the
positive cases of the formal theorem ensure existence in the stated
calculus; they need not provide a short explicit logarithmic formula
until a constructive five-term reduction is supplied. The consecutive
family and $3/5$ evaluation provide such explicit analytic formulas
in useful cases.

The adjacent-argument functional identity used for these evaluations
is an independently derived reformulation of the known Choie--Kumar
functional equation \cite[Theorem 3.2(1)]{ChoieKumar2023}. It should be
credited accordingly. We do not claim priority for a general functional
equation already present in that paper. The inspected follow-up work
on character-weighted Herglotz integrals supplies relevant future
directions \cite{DixitSathyanarayanaSharan2024}.

\subsection{Status synchronization across incoming reports}

The recent rigidity report already proves the $S_4$ identity by an
explicit combination of standard rows and a Möbius-duality row
\cite{PriorRigidity}. Older descriptions of $S_4$ as an unresolved
global identity should be updated when that report is integrated.
The earlier obstruction for a smaller relation space remains a valid
statement about that smaller space. It should be preserved with its
scope, not treated as a contradiction to the later proof.

Likewise, the residual level-three and level-four invariant spaces
were computed in the preceding report \cite{PriorResearch}. The
new contribution is the free same-point lift and the full ranks under
the manuscript's exact boundary convention. The supplied integration
guide records this dependency to avoid duplicate priority claims.
